% SPDX-License-Identifier: GPL-3.0-or-later
% Copyright (C) 2026 Alessandro Greco
%% Русское альтернативное описание EPUB.
\altfig{Пять криптографических задач Tor: конфиденциальность с AES-CTR для ячеек и AES-GCM либо ChaCha20-Poly1305 для каналов TLS (главы 10 и 7); целостность с дайджестами SHA-1 и SHA3-256 и тегами CGO на основе POLYVAL (глава 10); согласование ключей X25519 в ntor и ntor-v3 (глава 9); аутентификация Ed25519 и устаревающим RSA в сертификатах и консенсусе (главы 7 и 12); вывод ключей HKDF-SHA256 и SHAKE-256 (глава 9).}
% Пять криптографических задач Tor и используемые примитивы.
\begin{tikzpicture}[font=\sffamily\scriptsize,
  compito/.style={draw=torviola,fill=torviola!8,rounded corners=2pt,minimum width=2.6cm,minimum height=0.7cm},
  prim/.style={draw=torverde!60!black,fill=torverde!12,rounded corners=2pt,text width=3.9cm,align=center,minimum height=0.7cm},
  dove/.style={font=\sffamily\tiny,text=torgrigio,anchor=west}]
  \foreach \i/\c/\p/\d in {%
    0/Конфиденциальность/{AES-CTR; TLS с AES-GCM или ChaCha20-Poly1305}/{ячейки (гл.~10), каналы (гл.~7)},
    1/Целостность/{SHA-1 и SHA3-256, POLYVAL}/{дайджесты ячеек RELAY, теги аутентификации CGO (гл.~10)},
    2/{Согласование ключей}/{X25519 (Diffie-Hellman)}/{ntor, ntor-v3 (гл.~9)},
    3/Аутентификация/{Ed25519, RSA (постепенно выводится)}/{сертификаты, консенсус (гл.~7, 12)},
    4/Вывод ключей/{HKDF-SHA256, SHAKE-256}/{ключи цепи (гл.~9)}} {
    \node[compito] (c\i) at (0,-\i*1.0) {\c};
    \node[prim] (p\i) at (4.4,-\i*1.0) {\p};
    \draw[->,torgrigio] (c\i) -- (p\i);
    \node[dove] at (6.8,-\i*1.0) {\d};
  }
\end{tikzpicture}
